\chapter*{序}
\addcontentsline{toc}{chapter}{序}
\markboth{序}{序}

从初中进入高中，物理学习最明显的变化，不只是公式变多、题目变长，而是描述问题的方式发生了变化。生活语言需要被翻译成物理量，直观情境需要被简化成模型，结论还要说明适用条件。许多学生在这个阶段感到困难，往往并不是缺少计算能力，而是还没有建立起从情境、图像到规律的完整通路。

编写这本讲义，是希望为这段过渡提供一条可以循序前进的路线。书名中的“一个月打通”，不是承诺三十天内掌握高中物理的全部难题，也不是用三十组公式替代完整课程。这里的“打通”，指的是先打通几条最重要的连接：从生活经验到物理定义，从文字条件到规范图示，从单个公式到适用条件，从局部计算到结果复核。只要这些连接逐步建立，后续学习就不再只是记忆零散结论。

\section*{一、为什么按三十节编排}

全书分为四部分。Day 01--07 从物理量、运动描述和受力分析起步；Day 08--14 进入功和能、动量、圆周运动、万有引力、振动与波；Day 15--24 建立电场、电路、磁场和电磁感应的主干；Day 25--30 扩展到光学、热学、近代物理、实验方法、模型复盘和综合训练。

这样的顺序遵循三个原则：先有语言再有方程，先有直观图像再有抽象关系，先能处理单一模型再练习模型切换。每一部分前的“学习路线”页说明本部分的知识台阶和困难回退位置；Day 29、Day 30 则把前面分散学习的模型重新连接起来。

\section*{二、这本讲义怎样讲物理}

本书把“对象、过程、条件、规律、结果”作为贯穿始终的五个问题。面对一道题，先确认研究对象和过程边界，再画出运动、受力、场、电路或光路图；使用公式时说明条件，得到结果后检查方向、单位、数量级和极限情况。

图像在本书中不是装饰。受力图用来确认力的来源和方向，$x$-$t$ 与 $v$-$t$ 图用来表达运动关系，电路图用来呈现连接和能量转化，场线、轨迹、光路、$p$-$V$ 图和能级图则承担各自的物理信息。读者应当尝试从图中说出对象、方向、几何关系和变化过程，而不是只看图形是否熟悉。

例题和变式训练也不以追求偏、难、怪为目标。例题用于展示模型选择和推理过程，变式训练用于检验同一方法能否迁移。第八板块的拓展内容提供竞赛入门或更深入的学科视角；它建立在本节主线已经理解的基础上，不应替代基本概念和典型模型的学习。

\section*{三、写给学生、教师和家长}

对学生来说，学习速度不必完全一致。看不懂时，应先回到图、定义和条件，而不是立即增加刷题数量。能把一道题为什么这样画、为什么这样列式、结果为什么可信说清楚，比只得到一个数值更重要。

对教师和家长来说，这本讲义更适合作为预习、衔接和复习材料，而不是替代课堂教学。读者的知识基础、课程版本和学习节奏各不相同，可依据“怎样使用这本讲义”中的分层建议选择内容。基础主线尚未稳定时，可以暂缓拓展部分；学有余力时，再回看其中的推导、模型边界和科学发展背景。

物理学不断发展，但物理学习中尊重事实、明确条件、建立模型和检验结论的方法具有长久价值。希望读者完成这三十节后，不只是多记住一些公式，也能更有把握地观察问题、提出假设、画出结构、作出推理，并愿意用证据修正自己的判断。

\begin{flushright}
董兴习\quad 许琴\\
2026 年 7 月
\end{flushright}
